
\chapter{Introducción}
\label{chap:introduccion}

\section{Planteamiento del problema}
\label{sec:planteamiento}

La volatilidad implícita es un concepto central en mercados de derivados. En una opción negociada, el precio observable resume expectativas, oferta y demanda, liquidez, microestructura y percepción de riesgo. Sin embargo, para comparar contratos con strikes y vencimientos distintos, los profesionales no suelen razonar únicamente en precios: transforman esos precios en volatilidades implícitas y analizan smiles, skews y estructuras temporales \citep{hull2022options,derman2016volatility}.

El problema de este trabajo surge de una doble exigencia en finanzas cuantitativas. Por un lado, los modelos de aprendizaje automático pueden aproximar relaciones no lineales entre strike, subyacente, vencimiento, tipo de interés y volatilidad. Por otro lado, en un entorno financiero no basta con obtener una predicción numéricamente competitiva. El modelo debe poder auditarse: hay que saber dónde falla, qué variables explican sus predicciones, si respeta patrones razonables de la superficie y si una predicción concreta está respaldada por datos históricos similares. Esta exigencia es coherente con la literatura reciente sobre explicabilidad y gobierno de modelos en finanzas, que enfatiza transparencia, interpretabilidad, documentación y revisión independiente cuando se usan modelos complejos en procesos críticos \citep{solaimani2025beyond,perezcruz2025managing}.

En este contexto, el proyecto construye modelos de volatilidad explicable, consultable y útil vía API. La explicabilidad se entiende en dos niveles. A nivel de ingeniería, sirve para detectar sesgos, errores de datos, incoherencias, discontinuidades y comportamientos no deseados. A nivel financiero y profesional, permite justificar decisiones ante revisiones internas, auditorías, model validation o análisis de riesgo.

\textbf{Operativamente, el problema se concreta en cuatro preguntas.}

\begin{itemize}
  \item Cómo transformar operaciones reales de mercado en una base de volatilidad implícita que sea consistente, trazable y reutilizable.
  \item Cómo aproximar la superficie \(\sigma(K,T)\) con modelos suficientemente flexibles sin introducir sesgos temporales o dependencias espurias.
  \item Cómo operacionalizar ese modelo de forma utilizable, de manera que no quede encerrado en notebooks o tablas estáticas.
  \item Cómo saber cuándo una predicción es robusta y cuándo conviene tratarla con cautela porque está en una región poco soportada o mal explicada.
\end{itemize}

La memoria se construye alrededor de esa cadena completa. No se limita a entrenar un predictor, sino que conecta datos de mercado, inversión de Black-76, selección de familias, explicabilidad global y local, diagnóstico por zonas de la superficie y despliegue. Ese enfoque es coherente con el tipo de uso que tendría el sistema en un entorno profesional: consulta repetida, revisión crítica y necesidad de justificar el resultado más allá de una métrica agregada. La construcción detallada de datos y ETL se amplía en el apéndice \ref{app:data-structures}; la parte de explicabilidad se desarrolla con más detalle en los apéndices \ref{app:shap}, apéndice \ref{app:ale} y apéndice \ref{app:symbolic-regression}; y la dimensión de despliegue y trazabilidad de ingeniería se complementa en el apéndice \ref{app:workflow-dev}.

Estas cuatro preguntas se responden de forma progresiva en el \autoref{chap:metodologia} (construcción y validación del pipeline), en el \autoref{chap:resultados} (evidencia empírica y lectura operativa) y en el \autoref{chap:conclusiones} (síntesis final y alcance práctico).

\section{Relevancia financiera y profesional}
\label{sec:relevancia}

La volatilidad implícita condensa información de mercado y se utiliza en valoración, cobertura, gestión de riesgos y análisis de escenarios. En opciones sobre índices, la forma de la superficie de volatilidad puede reflejar demanda de protección, asimetría percibida, presión de liquidez o cambios de régimen. Un modelo que suaviza en exceso los extremos de moneyness, que comprime volatilidades extremas o que presenta saltos artificiales en vencimientos cortos puede producir señales engañosas aunque su RMSE agregado sea aceptable.

La literatura clásica de derivados establece que la volatilidad no debe interpretarse como un simple parámetro estadístico fijo, sino como una magnitud inferida del precio bajo un modelo de valoración \citep{black1976pricing,hull2022options}. La literatura específica sobre smile muestra que los mercados organizan la volatilidad en torno a moneyness y vencimiento, no solo en torno a niveles absolutos de strike \citep{derman2016volatility}. Por ello, el proyecto se diseña desde la geometría financiera de la superficie: el aprendizaje automático se usa para aproximar una función de volatilidad, pero la evaluación exige interpretar su forma.

\textbf{La relevancia práctica aparece en varios frentes.}

\begin{itemize}
  \item \textbf{Valoración y marking:} una estimación coherente de volatilidad implícita permite comparar contratos y revisar precios fuera de rango.
  \item \textbf{Gestión de riesgos:} la sensibilidad del error por moneyness y vencimiento importa tanto como el error medio total.
  \item \textbf{Gobierno de modelos:} el área usuaria necesita entender si el modelo reacciona con lógica financiera y no solo si ajusta bien una muestra histórica.
  \item \textbf{Infraestructura analítica:} si el modelo no puede consultarse, versionarse y explicarse, su valor operativo cae de forma drástica.
\end{itemize}

Desde esa perspectiva, un buen resultado no es únicamente ``predecir bien''. También implica poder distinguir entre una predicción bien soportada, una extrapolación razonable y una salida que requiere revisión adicional. Ese matiz es el que justifica que la memoria dedique espacio tanto a finanzas como a IA y a arquitectura de servicio.

\section{Antecedentes}
\label{sec:antecedentes}

La valoración de opciones sobre futuros se apoya tradicionalmente en Black-76 \citep{black1976pricing}. En ese marco, el precio de una opción europea depende del futuro subyacente, el strike, el vencimiento, el tipo de interés y la volatilidad. Cuando el precio de mercado se observa, la volatilidad implícita se obtiene resolviendo numéricamente la ecuación inversa. En la literatura reciente, las redes neuronales se han utilizado tanto para acelerar la valoración y el cálculo de volatilidades implícitas \citep{liu2019pricing} como para estudiar movimientos de la superficie de volatilidad en función de moneyness, vencimiento y entorno de volatilidad \citep{cao2019neural}.

En paralelo, el uso de modelos de aprendizaje automático para superficies financieras responde a la necesidad de capturar no linealidades, interacciones y patrones empíricos que no quedan completamente descritos por modelos paramétricos simples. Sin embargo, la adopción profesional de modelos flexibles exige herramientas de interpretación. SHAP proporciona una semántica aditiva para atribuir predicciones a variables \citep{lundberg2017unified}; ICE y ALE ayudan a estudiar respuestas funcionales \citep{goldstein2015peeking,apley2020visualizing}; los árboles surrogate y la regresión simbólica aproximan modelos complejos mediante objetos interpretables \citep{cranmer2023interpretable}.

\textbf{Los antecedentes relevantes convergen en tres líneas.}

\begin{itemize}
  \item \textbf{Línea financiera:} la inversión de modelos como Black-76 convierte precios observados en volatilidades comparables y lleva el problema al terreno de la superficie.
  \item \textbf{Línea empírica:} los smiles y las estructuras temporales muestran que la volatilidad de mercado no es constante y que la relación con strike y vencimiento es no lineal.
  \item \textbf{Línea metodológica:} los modelos de caja negra ganan flexibilidad, pero obligan a complementar la predicción con mecanismos de explicación y validación.
\end{itemize}

Este trabajo se sitúa precisamente en la intersección de esas tres líneas. Toma la formulación clásica de derivados como base para construir la variable objetivo, utiliza familias de aprendizaje automático sobre datos tabulares para aproximar la función de volatilidad y añade una capa de explicabilidad y servicio para hacer la solución defendible en un contexto profesional.

\section{Objetivos}
\label{sec:objetivos}

El objetivo general es desarrollar una solución reproducible para estimar, explicar y consultar volatilidad implícita de opciones del IBEX. Este objetivo se concreta en los siguientes objetivos específicos:

\begin{enumerate}
  \item Construir una base de volatilidad implícita a partir de operaciones, contratos, futuros subyacentes y curvas de tipos, con validaciones de calidad y trazabilidad.
  \item Diseñar un conjunto de variables financieras coherentes con la geometría de la superficie de volatilidad.
  \item Entrenar y comparar varias familias de modelos de regresión bajo una metodología temporal que reduzca lookahead bias y data snooping.
  \item Generar componentes de explicabilidad global, local, funcional y de diagnóstico.
  \item Implementar una API y un dashboard que permitan consultar predicciones, explicar muestras concretas y revisar el comportamiento del modelo.
  \item Preparar infraestructura cloud reproducible para almacenar modelos, paquetes explicativos y metadatos de servicio.
\end{enumerate}

Tomados en conjunto, estos objetivos cubren los tres planos que estructuran la memoria. El primero es \textbf{financiero}, porque la variable objetivo y las transformaciones deben respetar la lógica de opciones sobre futuros. El segundo es \textbf{de IA y Data Science}, porque la comparación entre familias exige un diseño temporal riguroso, métricas correctas y lectura crítica del error. El tercero es \textbf{de ingeniería}, porque el resultado final no es un gráfico aislado, sino un sistema consultable y versionable.

\section{Contribución del trabajo}
\label{sec:contribucion}

La principal contribución del trabajo es integrar en un único flujo operativo elementos que habitualmente se tratan por separado: construcción de datos de mercado, inversión de Black-76, entrenamiento de modelos, explicabilidad, diagnóstico financiero, API y despliegue cloud. La Figura~\ref{fig:doc-index-1} resume esta cadena con las mismas dependencias que la documentación técnica.

\DiagramIndexOne

La aportación académica se sitúa en la metodología de modelización y validación; la aportación profesional, en la capacidad de convertir un modelo de volatilidad en una herramienta inspeccionable y desplegable.

\textbf{De forma más concreta, la contribución del trabajo puede resumirse así:}

\begin{itemize}
  \item \textbf{Contribución financiera:} construir una base de volatilidad implícita a partir de datos reales enlazando operaciones de opciones, futuros y curvas de tipos.
  \item \textbf{Contribución metodológica:} comparar familias heterogéneas con el mismo protocolo temporal, las mismas métricas y el mismo criterio de selección.
  \item \textbf{Contribución de explicabilidad:} combinar atribuciones SHAP, curvas ICE y ALE, modelos equivalentes, superficies y vecinos históricos en un único flujo de lectura.
  \item \textbf{Contribución de ingeniería:} dejar el modelo convertido en un paquete autocontenido consumible por dashboard y API, con almacenamiento reproducible y despliegue desacoplado.
\end{itemize}

\textbf{Desde el punto de vista de negocio en derivados, el valor aplicado se concreta en ocho usos directos:}

\begin{itemize}
  \item \textbf{Valoración de opciones ilíquidas o poco negociadas:} la predicción se apoya en estructura de superficie y soporte histórico de vecinos, no solo en una cotización puntual.
  \item \textbf{Interpolación y suavizado de superficies de volatilidad:} el dashboard permite revisar continuidad por moneyness y vencimiento y detectar escalones no deseados.
  \item \textbf{Detección de precios anómalos o errores de mercado:} los residuales por zonas de superficie ayudan a localizar observaciones fuera de patrón.
  \item \textbf{Validación independiente de modelos internos:} el pipeline puede usarse como contraste externo frente a pricers de mesa o librerías internas.
  \item \textbf{Control de riesgos de mesa de derivados:} el error deja de leerse solo en promedio y se monitoriza por tramos relevantes de moneyness y vencimiento.
  \item \textbf{Explicación ante auditoría interna o model validation:} cada predicción puede justificarse con SHAP, curvas funcionales y evidencia de soporte histórico.
  \item \textbf{Análisis de sensibilidad por moneyness y vencimiento:} superficies locales, ICE y ALE aportan lectura funcional de cómo reacciona el modelo.
  \item \textbf{Soporte a decisiones de cobertura:} la lectura conjunta de smile, term structure y diagnóstico local ofrece contexto operativo para coberturas.
\end{itemize}

Esa combinación es la que hace que el proyecto no sea solo un ejercicio de entrenamiento de modelos. El valor añadido está en que cada capa técnica responde a una pregunta concreta y útil: cómo se construye la volatilidad, cómo se estima, cómo se valida, cómo se explica y cómo se despliega.

\section{Estructura de la memoria}
\label{sec:estructura}

El \autoref{chap:marco} presenta los fundamentos financieros, de aprendizaje automático, explicabilidad e infraestructura. El \autoref{chap:metodologia} describe la construcción de datos, variables, modelos, validación, dashboard, API y cloud. El \autoref{chap:resultados} analiza las métricas persistidas y la lectura que debe hacerse desde las vistas del dashboard. El \autoref{chap:conclusiones} resume conclusiones y líneas futuras. Los apéndices desarrollan con mayor profundidad la estructura de datos y ETL (apéndice \ref{app:data-structures}), la metodología de desarrollo y despliegue (apéndice \ref{app:workflow-dev}), SHAP (apéndice \ref{app:shap}), ALE (apéndice \ref{app:ale}), el modelo Quantum Inspired (apéndice \ref{app:quantum-inspired}), la regresión simbólica (apéndice \ref{app:symbolic-regression}), el entrenamiento progressive ATM (apéndice \ref{app:progressive}) y el análisis de sensibilidad por moneyness y vencimiento (apéndice \ref{app:sensibilidad-moneyness-vencimiento}).

\textbf{El orden de lectura responde al flujo real del proyecto.}

\begin{enumerate}
  \item En la \textbf{Introducción} se delimita el problema y se justifica por qué la explicabilidad es parte del resultado, no un añadido posterior.
  \item En el \textbf{Marco Teórico} se fijan las piezas conceptuales necesarias para interpretar Black-76, smiles, familias de modelos, métricas y gobierno de modelos.
  \item En la \textbf{Metodología} se documenta la implementación efectiva: ETL, features, validación temporal, reentrenamiento, generación del paquete de dashboard y despliegue.
  \item En \textbf{Resultados y Discusión} se conectan métricas, superficies, explicaciones y limitaciones con una lectura financiera y técnica unificada.
  \item En los \textbf{Apéndices} se amplían los mecanismos de explicabilidad y entrenamiento para que la memoria principal conserve foco sin perder profundidad técnica.
\end{enumerate}

